\scanpage{151}
Si $2$ est inv., il y a alors un seul pt fixe dans $A$, et $A$ devient ainsi un module, où l'\unsure{endom|ordinaire} n'est autre que $x\mapsto-x$. Si par contre $2$ est \deleted{inversibl} \inserted{nul} ds $k$, \deleted{\unsure{l'est}} l'élément $t$ est dans le module des translations, et la donnée équivaut à celle de $A$, et de $t\in A^0$ --- et l'involution sera la translation par $t$. Nous nous intéressons cependant aux cas ``à cheval'', tel celui où $k=\mathbb Z$, quand $2$ n'est ni inv., ni nul dans $k$. \deleted{Je veux définir}

Je voudrais définir une opération de produit tensoriel d'espaces affines à involution. \unsure{Soit}\deleted{\unsure{ent}} $(A_1,A_2)$ deux tels espaces, avec involutions $\sigma_1,\sigma_2$ (je laisse tomber la notation $x'$, qui est \unsure{ambigue|ambiguë} \unsure{qqfois|parfois}) et soit $B$ un espace affine sur $k$, les espaces de translation sont notés $A_1^0,A_2^0,B^0$.
% slanted
\margin{il suffit que $B$ soit un \uline{ens.} sur lequel un \deleted{\unsure{groupe}} module $B^0$ opère (\unsure{sans} pt fixe \unsure{et transitif.})}
Une application
\[
 f:A_1\times A_2\longrightarrow B
\]
est dite \deleted{\unsure{une}} $(\sigma_1,\sigma_2)$-bilinéaire, si \deleted{$\unsure{\exists!|\exists}$}
\[
 \deleted{f(x}\quad\exists f^0:A_1^0\times A_2^0\longrightarrow B^0\qquad\text{bil}
\]
% archive sheet number: 150

